\section{Limitations}
\label{sec:limitations}

% Keep this section concise. Write it as four connected paragraphs, not as a
% separate warning for every technical decision. The fuller checklist remains
% in the project notes for fact-checking while drafting.

\TODO{PARAGRAPH 1 --- WHAT THE MODEL CAN AND CANNOT SHOW: Explain that this is a
model of how connected the landscape appears, not proof of where cheetahs actually
moved. The cores came from 2010--2016 records and were kept in the same places for
every year. This makes the comparison fair, but it means the analysis measures change
around the cores rather than change in cheetah populations. Also mention that the
starting records may be uneven because some areas were surveyed more than others.}

\TODO{PARAGRAPH 2 --- ASSUMPTIONS IN THE MODEL: Explain that the resistance values
were chosen using published research and clearly stated assumptions because suitable
GPS movement data were not available. The sensitivity tests show whether the main
results survive different reasonable assumptions, but they cannot prove that one set
of values exactly represents cheetah movement. Include the important vegetation
detail in one sentence: this part of the model mostly measures bare ground rather than
all aspects of vegetation structure.}

\TODO{PARAGRAPH 3 --- LIMITS OF THE INPUT DATA: Explain that roads, livestock,
terrain, and fences shaped the map but were held constant through time, so they could
not create year-to-year change. Note that data completeness may differ among
countries, rainfall was not modeled, the nighttime-light product changed version in
2024, and the 2012 built-up snapshot used 2010 data. These issues mean that some small
changes may reflect the source data or short-term environmental variation rather than
lasting landscape change. Do not list every dataset again; keep these examples in one
paragraph.}

\TODO{PARAGRAPH 4 --- SCALE AND LOCATION UNCERTAINTY: Explain that the study used
one core definition and a 1-km grid, so results may differ at another scale. The model
was more consistent about the relative cost of connecting two cores than about the
exact line a route should follow: 49 of 180 route-and-year comparisons had less than
80 percent overlap within 1 km. Therefore, treat the mapped routes as broad areas for
planning and field checking, not exact travel lines. Briefly state that results should
not be transferred to another region without testing them there. If space permits,
note that current was hidden inside cores for clearer maps and that the resistance
grid contained 34 disconnected pieces; otherwise place those two technical details in
the supplement.}
